% 第4章第2节（下）：矩阵运算的性质，按运算分组（加法与数乘、转置、乘法、逐项乘积、逆、幂、迹、行列式、范数），每条性质先给结论再给证明或说明；常见矩阵（含正交矩阵）的性质分别放在相应运算之下。

\subsection{矩阵运算的性质}
\label{sec:ls-matrix-operation-properties}
\label{sec:prelim-matrix-operation-rules}
\label{sec:linear-algebra-map-composition}

上一小节只给出了各运算的算法；推导时还要知道哪些变形允许做。实数乘法的哪些规则保留下来，哪些不再成立？转置、求逆与乘法相遇时次序怎样变化？迹与行列式在乘法下怎样表现？矩阵范数怎样控制连续两次作用的放大？对角、三角、置换与正交矩阵在这些运算下又有哪些简单规律？本小节按运算分组回答，每条性质先写出结论，再给出证明或说明。以下凡写出乘积，都默认各因子的形状相容。

\subsubsection{加法与数乘的性质}

同形矩阵逐项相加与数乘，向量一节的八条规则在每个位置上照搬过来。对同形矩阵$\bm A,\bm B,\bm C$和实数$a,b$：
\[
\begin{aligned}
 &\bm A+\bm B=\bm B+\bm A, && (\bm A+\bm B)+\bm C=\bm A+(\bm B+\bm C),\\
 &\bm A+\bm0=\bm A, && \bm A+(-\bm A)=\bm0,\\
 &a(b\bm A)=(ab)\bm A, && 1\bm A=\bm A,\\
 &(a+b)\bm A=a\bm A+b\bm A, && a(\bm A+\bm B)=a\bm A+a\bm B.
\end{aligned}
\]
每一条都可以在任一位置$(i,j)$上用实数运算直接验证。

\subsubsection{转置的性质}
\label{sec:ls-reshaping}

\paragraph{与加法、数乘交换}
\[
 (\bm A+\bm B)^\top=\bm A^\top+\bm B^\top,\qquad
 (c\bm A)^\top=c\bm A^\top,\qquad
 (\bm A^\top)^\top=\bm A.
\]
转置只搬动元素，逐项运算的结果搬到新位置后不变；转置两次，每个元素回到原来的位置。

\paragraph{乘积的转置}
\[
 (\bm A\bm B)^\top=\bm B^\top\bm A^\top.
\]
乘积转置后因子次序反转。按元素公式，
\[
 \bigl((\bm A\bm B)^\top\bigr)_{i,j}=(\bm A\bm B)_{j,i}=\sum_\ell a_{j,\ell}b_{\ell,i}
 =\sum_\ell(\bm B^\top)_{i,\ell}(\bm A^\top)_{\ell,j}=(\bm B^\top\bm A^\top)_{i,j}.
\]
多个因子时依次使用，$(\bm A\bm B\bm C)^\top=\bm C^\top\bm B^\top\bm A^\top$。

\paragraph{内积中的转置}
\[
 \langle\bm A\bm x,\bm y\rangle=\langle\bm x,\bm A^\top\bm y\rangle,\qquad\bm A\in\mathbb R^{m\times d},\ \bm x\in\mathbb R^d,\ \bm y\in\mathbb R^m.
\]
由上一条，$(\bm A\bm x)^\top\bm y=\bm x^\top\bm A^\top\bm y$：在欧氏内积中，作用在左侧向量上的$\bm A$可以改写成作用在右侧向量上的$\bm A^\top$。

\paragraph{$\bm A^\top\bm A$对称}
对任意矩阵$\bm A$，$\bm A^\top\bm A$与$\bm A\bm A^\top$都是对称矩阵，因为$(\bm A^\top\bm A)^\top=\bm A^\top(\bm A^\top)^\top=\bm A^\top\bm A$，另一个同理。下面直接算出一个$\bm A\bm A^\top$。
\begin{example}[计算$\bm A\bm A^\top$]
\label{ex:ls-gram-product-symmetric}
给定：$\bm A=\left[\begin{smallmatrix}1&2&3\\4&5&6\end{smallmatrix}\right]$。
计算：$\bm A\bm A^\top$的$(i,j)$元是$\bm A$的第$i$行与第$j$行的内积，
\[
 \bm A\bm A^\top=\begin{bmatrix}1+4+9&4+10+18\\4+10+18&16+25+36\end{bmatrix}=\begin{bmatrix}14&32\\32&77\end{bmatrix}.
\]
结论：两个非对角元都是两行的内积$32$，结果是对称矩阵。
\end{example}

\paragraph{正交矩阵的转置条件}
$n$阶方阵$\bm Q$是正交矩阵，当且仅当
\[
 \bm Q^\top\bm Q=\bm I_n.
\]
$(\bm Q^\top\bm Q)_{i,j}$是$\bm Q$的第$i$列与第$j$列的内积。所以“各列都是单位向量”对应对角元等于$1$，“各列两两正交”对应非对角元等于$0$，两者合起来正是这个矩阵等式。下面用这个条件重新检验示例~\ref{ex:ls-rotation-orthogonal}中的两个矩阵。
\begin{example}[用转置条件检验正交矩阵]
\label{ex:ls-orthogonality-by-transpose}
给定：旋转矩阵$\bm R_\phi$与对角矩阵$\operatorname{diag}(2,1)$。
计算：
\[
 \begin{aligned}
 \bm R_\phi^\top\bm R_\phi
 &=\begin{bmatrix}\cos\phi&\sin\phi\\-\sin\phi&\cos\phi\end{bmatrix}
   \begin{bmatrix}\cos\phi&-\sin\phi\\\sin\phi&\cos\phi\end{bmatrix}\\
 &=\begin{bmatrix}\cos^2\phi+\sin^2\phi&0\\0&\sin^2\phi+\cos^2\phi\end{bmatrix}
 =\bm I_2,\\
 \operatorname{diag}(2,1)^\top\operatorname{diag}(2,1)&=\operatorname{diag}(4,1)\neq\bm I_2.
 \end{aligned}
\]
结论：$\bm R_\phi$满足转置条件，是正交矩阵；$\operatorname{diag}(2,1)$的乘积在$(1,1)$位置为$4\neq1$，对应第一列长度不为$1$。
\end{example}

\paragraph{正交矩阵保持内积}
\[
 \langle\bm Q\bm x,\bm Q\bm y\rangle=\langle\bm x,\bm y\rangle\qquad(\bm Q\text{为正交矩阵}).
\]
由内积中的转置与上一条，$\langle\bm Q\bm x,\bm Q\bm y\rangle=\langle\bm x,\bm Q^\top\bm Q\bm y\rangle=\langle\bm x,\bm y\rangle$。特别地，正交矩阵把两两正交的向量仍送到两两正交的向量。

\subsubsection{乘法的性质}

\paragraph{单位乘}
\[
 \bm I_m\bm A=\bm A=\bm A\bm I_d,\qquad\bm A\in\mathbb R^{m\times d}.
\]
$\bm I_m$的第$i$行是$\bm e_i^\top$，它与$\bm A$的第$j$列相乘取出$a_{i,j}$；右乘同理。单位矩阵在矩阵乘法中的作用相当于数$1$在实数乘法中的作用。

\paragraph{零乘}
\[
 \bm0\bm A=\bm0,\qquad\bm A\bm0=\bm0.
\]
乘积的每个元素都是若干乘积之和，其中一个因子全为零时，每一项都为零。这里各个零矩阵的形状由相乘的另一个因子决定。

反过来不成立：$\bm A\bm B=\bm0$不能推出$\bm A=\bm0$或$\bm B=\bm0$，下面是一个反例。
\begin{example}[反例：两个非零矩阵的乘积为零]
\label{ex:ls-zero-product-nonzero-factors}
给定：$\bm A=\left[\begin{smallmatrix}1&0\\0&0\end{smallmatrix}\right]$，$\bm B=\left[\begin{smallmatrix}0&0\\1&0\end{smallmatrix}\right]$，两者都不是零矩阵。
计算：
\[
 \bm A\bm B=\begin{bmatrix}1\cdot0+0\cdot1&1\cdot0+0\cdot0\\0\cdot0+0\cdot1&0\cdot0+0\cdot0\end{bmatrix}=\begin{bmatrix}0&0\\0&0\end{bmatrix}.
\]
结论：$\bm A\bm B=\bm0$，而$\bm A\neq\bm0$、$\bm B\neq\bm0$。
\end{example}
原因可以按列看出。由乘积的元素公式，$\bm A\bm B$的第$j$列就是$\bm A$乘$\bm B$的第$j$列，所以$\bm A\bm B=\bm0$恰好意味着$\bm B$的每一列$\bm b$都满足$\bm A\bm b=\bm0$：右因子产生的结果被左因子完全消去。上例中$\bm A$只保留第一个坐标、把第二个坐标清零，而$\bm B$的第一列$(0,1)^\top$只有第二个坐标非零，第二列为零。

\paragraph{结合律}
\[
 (\bm A\bm B)\bm C=\bm A(\bm B\bm C).
\]
按式~\eqref{eq:prelim-matrix-composition}展开，两边的第$(i,j)$项都是$\sum_k\sum_l a_{i,k}b_{k,l}c_{l,j}$，只是两个有限求和的次序不同。这也是函数复合满足结合律（第~\ref{sec:prelim-map-operations}节）在矩阵上的翻译。由结合律，多个相容矩阵连乘时不必标出括号。

\paragraph{分配律与数乘相容}
\[
 \bm A(\bm B+\bm C)=\bm A\bm B+\bm A\bm C,\qquad
 (\bm A+\bm B)\bm C=\bm A\bm C+\bm B\bm C,\qquad
 c(\bm A\bm B)=(c\bm A)\bm B=\bm A(c\bm B),
\]
逐项展开有限和即得。把右因子取为列向量，得到矩阵与向量的乘积对输入保持线性组合：
\begin{equation}
 \bm A(a\bm u+b\bm v)=a\bm A\bm u+b\bm A\bm v,
 \qquad\bm u,\bm v\in\mathbb R^d,\quad a,b\in\mathbb R.
 \label{eq:prelim-matrix-linearity}
\end{equation}

由这一性质，函数$\bm u\mapsto\bm A\bm u$由它在$\bm e_1,\ldots,\bm e_d$上的输出完全确定：因为$\bm u=\sum_ju_j\bm e_j$，所以$\bm A\bm u=\sum_ju_j\bm A\bm e_j$，而这些输出$\bm A\bm e_j$就是$\bm A$的各列。第~\ref{chap:linear-representation}章把保持线性组合的函数称为线性映射，并说明每个这样的映射都能由一张矩阵记录。

\paragraph{分块乘法}
把矩阵按行、列分成尺寸相容的块，乘法规则对块照样成立：
\[
 \begin{bmatrix}\bm A&\bm B\end{bmatrix}\begin{bmatrix}\bm C\\\bm D\end{bmatrix}=\bm A\bm C+\bm B\bm D,
 \qquad
 \begin{bmatrix}\bm A\\\bm B\end{bmatrix}\bm C=\begin{bmatrix}\bm A\bm C\\\bm B\bm C\end{bmatrix},
\]
\[
 \bm A\begin{bmatrix}\bm C&\bm D\end{bmatrix}=\begin{bmatrix}\bm A\bm C&\bm A\bm D\end{bmatrix},
\]
以及$2\times2$分块的乘积
\[
 \begin{bmatrix}\bm A_{1,1}&\bm A_{1,2}\\\bm A_{2,1}&\bm A_{2,2}\end{bmatrix}
 \begin{bmatrix}\bm B_{1,1}&\bm B_{1,2}\\\bm B_{2,1}&\bm B_{2,2}\end{bmatrix}
 =\begin{bmatrix}
  \bm A_{1,1}\bm B_{1,1}+\bm A_{1,2}\bm B_{2,1}&\bm A_{1,1}\bm B_{1,2}+\bm A_{1,2}\bm B_{2,2}\\
  \bm A_{2,1}\bm B_{1,1}+\bm A_{2,2}\bm B_{2,1}&\bm A_{2,1}\bm B_{1,2}+\bm A_{2,2}\bm B_{2,2}
 \end{bmatrix}.
\]
条件是每个相乘的块对中，左块的列数等于右块的行数。第一式成立，是因为式~\eqref{eq:prelim-matrix-composition}对中间索引$\ell$的求和可以分段进行：前一段$\ell$对应$\bm A$的列与$\bm C$的行，后一段对应$\bm B$的列与$\bm D$的行，两段之和逐项等于整段求和。后两式只是把结果的行或列分组，$2\times2$分块由这三式合成。块可以像数一样参与运算，只是块之间的次序不能交换。

把块分到最细，得到三种常用的特例。右因子按列分块，得$(\bm A\bm B)_{:,j}=\bm A\bm B_{:,j}$：$\bm A$分别乘$\bm B$的每一列。左因子按行分块，得$(\bm A\bm B)_{i,:}=\bm A_{i,:}\bm B$：结果的第$i$行是$\bm B$各行以$\bm A_{i,:}$的元素为系数的线性组合。

第三种特例把中间维数$d$分成长度为$1$的$d$段，这时第一式给出
\begin{equation}
 \bm A\bm B=\sum_{\ell=1}^d\bm A_{:,\ell}\bm B_{\ell,:},
 \label{eq:ls-product-outer-sum}
\end{equation}
其中一列$\bm A_{:,\ell}\in\mathbb R^{m\times1}$乘一行$\bm B_{\ell,:}\in\mathbb R^{1\times k}$，得到$(i,j)$元素为$a_{i,\ell}b_{\ell,j}$的$m\times k$矩阵。式~\eqref{eq:ls-product-outer-sum}是后面低秩近似的出发点。下面用示例~\ref{ex:ls-matrix-product-2x2}的乘积核对这个公式。
\begin{example}[把乘积写成列乘行之和]
\label{ex:ls-product-as-outer-sum}
给定：示例~\ref{ex:ls-matrix-product-2x2}中的$\bm A=\left(\begin{smallmatrix}1&1\\0&1\end{smallmatrix}\right)$与$\bm B=\left(\begin{smallmatrix}1&0\\0&2\end{smallmatrix}\right)$。
计算：$\bm A$的两列为$(1,0)^\top$与$(1,1)^\top$，$\bm B$的两行为$(1,0)$与$(0,2)$，
\[
 \bm A\bm B=\begin{pmatrix}1\\0\end{pmatrix}\begin{pmatrix}1&0\end{pmatrix}+\begin{pmatrix}1\\1\end{pmatrix}\begin{pmatrix}0&2\end{pmatrix}
 =\begin{pmatrix}1&0\\0&0\end{pmatrix}+\begin{pmatrix}0&2\\0&2\end{pmatrix}=\begin{pmatrix}1&2\\0&2\end{pmatrix}.
\]
结论：两个“列乘行”矩阵之和与逐项计算的乘积相同。
\end{example}

\paragraph{交换律不成立}
一般$\bm A\bm B\neq\bm B\bm A$。形状不同时，两个乘积甚至不一定都有定义；即使两个因子是同阶方阵，两种次序的乘积也可能不同，下面是一个反例。
\begin{example}[反例：先交换再保留与先保留再交换]
\label{ex:ls-noncommuting-pair}
给定：只保留第一个坐标的$\bm P=\left[\begin{smallmatrix}1&0\\0&0\end{smallmatrix}\right]$，交换两个坐标的$\bm S=\left[\begin{smallmatrix}0&1\\1&0\end{smallmatrix}\right]$。
计算：
\[
 \bm P\bm S=\begin{bmatrix}0&1\\0&0\end{bmatrix},\qquad
 \bm S\bm P=\begin{bmatrix}0&0\\1&0\end{bmatrix}.
\]
结论：$\bm P\bm S\neq\bm S\bm P$。作用在$\bm u\in\mathbb R^2$上，$\bm P\bm S\bm u=(u_2,0)^\top$是先交换、再保留第一个坐标，$\bm S\bm P\bm u=(0,u_1)^\top$是先保留、再交换，两种次序的结果不同。
\end{example}
书写$\bm A\bm B\bm x$时，作用顺序是从右向左：先$\bm B$，后$\bm A$。

\paragraph{消去律不成立}
$\bm A\bm B=\bm A\bm C$且$\bm A\neq\bm0$不能推出$\bm B=\bm C$。原因是$\bm A\bm B=\bm A\bm C$等价于$\bm A(\bm B-\bm C)=\bm0$，而乘积为零不能推出$\bm B-\bm C=\bm0$。下面用示例~\ref{ex:ls-zero-product-nonzero-factors}构造反例。
\begin{example}[反例：非零的左因子不能消去]
\label{ex:ls-cancellation-fails}
给定：示例~\ref{ex:ls-zero-product-nonzero-factors}中的$\bm A=\left[\begin{smallmatrix}1&0\\0&0\end{smallmatrix}\right]$与$\bm B=\left[\begin{smallmatrix}0&0\\1&0\end{smallmatrix}\right]$，以及$\bm C=\bm0$。
验证：$\bm A\bm B=\bm0=\bm A\bm C$，而$\bm B\neq\bm C$。
结论：虽然$\bm A\bm B=\bm A\bm C$且$\bm A\neq\bm0$，却不能消去$\bm A$得到$\bm B=\bm C$。
\end{example}
若$\bm A$是可逆方阵，在$\bm A\bm B=\bm A\bm C$两边左乘$\bm A^{-1}$即得$\bm B=\bm C$。所以左因子可逆时可以消去，一般的非零因子则不能。

常见矩阵参与乘法时，结果有固定模式。下面逐类给出。

\paragraph{对角矩阵的乘法}
对$\bm D=\operatorname{diag}(d_1,\ldots,d_n)$，
\[
 (\bm D\bm A)_{i,j}=d_i\bm A_{i,j},\qquad(\bm A\bm D)_{i,j}=\bm A_{i,j}d_j,
\]
\[
 \operatorname{diag}(d_1,\ldots,d_n)\operatorname{diag}(e_1,\ldots,e_n)=\operatorname{diag}(d_1e_1,\ldots,d_ne_n).
\]
左乘对角矩阵缩放各行，右乘缩放各列：乘积公式中只有$\ell=i$（左乘时）或$\ell=j$（右乘时）的一项留下。作用在向量上，$\bm D\bm u=(d_1u_1,\ldots,d_nu_n)^\top$是逐坐标缩放。两个对角矩阵相乘时对角元逐个相乘，所以对角矩阵之间可以交换次序。下面的例子比较左乘与右乘。
\begin{example}[左乘与右乘对角矩阵]
\label{ex:ls-diagonal-scaling}
给定：$\bm D=\operatorname{diag}(2,3)$，$\bm A=\left[\begin{smallmatrix}1&2\\3&4\end{smallmatrix}\right]$。
计算：
\[
 \begin{aligned}
 \bm D\bm A&=\begin{bmatrix}2\cdot1&2\cdot2\\3\cdot3&3\cdot4\end{bmatrix}=\begin{bmatrix}2&4\\9&12\end{bmatrix},\\
 \bm A\bm D&=\begin{bmatrix}1\cdot2&2\cdot3\\3\cdot2&4\cdot3\end{bmatrix}=\begin{bmatrix}2&6\\6&12\end{bmatrix}.
 \end{aligned}
\]
结论：左乘把第一、二行分别乘以$2$、$3$，右乘把第一、二列分别乘以$2$、$3$。两个结果不同，所以对角矩阵与一般矩阵相乘时仍要区分左右。
\end{example}

\paragraph{三角矩阵的乘法}
若$\bm A,\bm B$都是$n$阶上三角矩阵，则$\bm A\bm B$也是上三角矩阵，并且
\[
 (\bm A\bm B)_{i,i}=\bm A_{i,i}\bm B_{i,i}.
\]
在$(\bm A\bm B)_{i,j}=\sum_\ell\bm A_{i,\ell}\bm B_{\ell,j}$中，$\bm A_{i,\ell}$只在$\ell\geq i$时可能非零，$\bm B_{\ell,j}$只在$\ell\leq j$时可能非零，所以只有$i\leq\ell\leq j$的项可能非零：$i>j$时没有这样的$\ell$，乘积元素为零；$i=j$时只剩$\ell=i$一项。下三角矩阵同理：两个下三角矩阵的乘积仍是下三角矩阵，对角元逐个相乘。下面是一个二阶算例。
\begin{example}[两个上三角矩阵的乘积]
\label{ex:ls-upper-triangular-product}
给定：$\bm A=\left[\begin{smallmatrix}1&2\\0&3\end{smallmatrix}\right]$，$\bm B=\left[\begin{smallmatrix}4&5\\0&6\end{smallmatrix}\right]$。
计算：
\[
 \bm A\bm B=\begin{bmatrix}1\cdot4+2\cdot0&1\cdot5+2\cdot6\\0\cdot4+3\cdot0&0\cdot5+3\cdot6\end{bmatrix}=\begin{bmatrix}4&17\\0&18\end{bmatrix}.
\]
结论：乘积仍是上三角矩阵，对角元$4=1\cdot4$、$18=3\cdot6$是两个因子对角元的乘积；右上角的$17$则同时用到了两个因子的非对角元。
\end{example}

\paragraph{置换矩阵的乘法}
设$\bm P=\left[\begin{smallmatrix}\bm e_{\sigma(1)}&\cdots&\bm e_{\sigma(n)}\end{smallmatrix}\right]$，其中$\sigma$是$\{1,\ldots,n\}$上的排列，则
\[
 (\bm P\bm A)_{\sigma(i),:}=\bm A_{i,:},\qquad(\bm A\bm P)_{:,j}=\bm A_{:,\sigma(j)}.
\]
即左乘把$\bm A$的第$i$行移到第$\sigma(i)$行，右乘使$\bm A\bm P$的第$j$列取自$\bm A$的第$\sigma(j)$列。例如$\left[\begin{smallmatrix}0&1\\1&0\end{smallmatrix}\right]$左乘交换两行、右乘交换两列。两个置换矩阵的乘积仍是置换矩阵：设$\bm P'$是另一个置换矩阵，$\bm P'\bm P$的第$j$列是$\bm P'\bm e_{\sigma(j)}$，即$\bm P'$的第$\sigma(j)$列；它是标准坐标向量，而且不同的$j$取到$\bm P'$的不同列，所以这些列互不相同。

\paragraph{正交矩阵的乘法}
若$\bm Q_1,\bm Q_2$是$n$阶正交矩阵，则
\[
 (\bm Q_1\bm Q_2)^\top(\bm Q_1\bm Q_2)=\bm Q_2^\top\bm Q_1^\top\bm Q_1\bm Q_2=\bm Q_2^\top\bm Q_2=\bm I_n,
\]
由正交矩阵的转置条件，$\bm Q_1\bm Q_2$仍是正交矩阵。下面的例子显示，两个平面旋转的乘积不仅是正交矩阵，而且仍是旋转。
\begin{example}[两个平面旋转的乘积]
\label{ex:ls-rotation-composition}
给定：旋转矩阵$\bm R_\phi$与$\bm R_\psi$。
计算：直接相乘，再用和角公式$\cos(\phi+\psi)=\cos\phi\cos\psi-\sin\phi\sin\psi$与$\sin(\phi+\psi)=\sin\phi\cos\psi+\cos\phi\sin\psi$，
\[
 \begin{aligned}
 \bm R_\phi\bm R_\psi
 &=\begin{bmatrix}\cos\phi\cos\psi-\sin\phi\sin\psi&-\cos\phi\sin\psi-\sin\phi\cos\psi\\\sin\phi\cos\psi+\cos\phi\sin\psi&-\sin\phi\sin\psi+\cos\phi\cos\psi\end{bmatrix}\\
 &=\bm R_{\phi+\psi}.
 \end{aligned}
\]
结论：先转$\psi$角、再转$\phi$角，等于一次转$\phi+\psi$角；乘积是正交矩阵，与上面的一般结论一致。
\end{example}

\paragraph{对称矩阵的乘积}
若$\bm A,\bm B$是同阶对称矩阵，则
\[
 (\bm A\bm B)^\top=\bm B^\top\bm A^\top=\bm B\bm A,
\]
所以$\bm A\bm B$对称当且仅当$\bm A\bm B=\bm B\bm A$。对称矩阵的乘积一般不对称，下面是一个反例。
\begin{example}[反例：两个对称矩阵的乘积不对称]
\label{ex:ls-symmetric-product-not-symmetric}
给定：对称矩阵$\bm A=\left[\begin{smallmatrix}1&1\\1&0\end{smallmatrix}\right]$与$\bm B=\operatorname{diag}(1,2)$。
计算：右乘对角矩阵缩放各列，
\[
 \bm A\bm B=\begin{bmatrix}1\cdot1&1\cdot2\\1\cdot1&0\cdot2\end{bmatrix}=\begin{bmatrix}1&2\\1&0\end{bmatrix}.
\]
结论：$(\bm A\bm B)_{1,2}=2\neq1=(\bm A\bm B)_{2,1}$，乘积不对称。这与上面的判据一致：$\bm B\bm A=\left[\begin{smallmatrix}1&1\\2&0\end{smallmatrix}\right]\neq\bm A\bm B$。
\end{example}

\subsubsection{逐项乘积的性质}

\paragraph{交换律}
\[
 \bm A\odot\bm B=\bm B\odot\bm A.
\]
式~\eqref{eq:ls-hadamard}右端只是实数乘法，按位置逐项验证即得。逐项乘积可交换而矩阵乘积不可交换，再次说明两者是不同的运算。

\paragraph{结合律}
\[
 (\bm A\odot\bm B)\odot\bm C=\bm A\odot(\bm B\odot\bm C),
\]
两边的$(i,j)$元素都是$a_{i,j}b_{i,j}c_{i,j}$。

\paragraph{分配律}
\[
 \bm A\odot(\bm B+\bm C)=\bm A\odot\bm B+\bm A\odot\bm C,
\]
两边的$(i,j)$元素都是$a_{i,j}(b_{i,j}+c_{i,j})$。

\paragraph{单位元是全$1$数组}
全部元素为$1$的同形数组$\bm J$满足$\bm J\odot\bm A=\bm A$。单位矩阵不是逐项乘积的单位元：$\bm I_2\odot\bm A$只保留$\bm A$的对角元，把非对角元都变为零。

\paragraph{对每个因子分别线性}
固定一个因子时，$\bm X\mapsto\bm M\odot\bm X$保持线性组合。两个因子都作为未知量时，乘积只对各自分别线性：同时把两者放大两倍，结果放大四倍。第~\ref{chap:linear-representation}章换一组记录方向以后会看到，矩阵乘积对应的复合关系保持不变，逐项乘积则随记录方式改变。

\subsubsection{逆的性质}

设下面的$\bm A,\bm B$都是$n$阶可逆方阵。

\paragraph{乘积的逆}
\begin{equation}
 (\bm A\bm B)^{-1}=\bm B^{-1}\bm A^{-1}.
 \label{eq:ls-inverse-product}
\end{equation}
先作用的后撤销，就像先穿袜子再穿鞋，脱的时候要先脱鞋。验证：$(\bm A\bm B)(\bm B^{-1}\bm A^{-1})=\bm A(\bm B\bm B^{-1})\bm A^{-1}=\bm A\bm A^{-1}=\bm I_n$，反方向$(\bm B^{-1}\bm A^{-1})(\bm A\bm B)=\bm I_n$同理。

\paragraph{逆的逆}
\[
 (\bm A^{-1})^{-1}=\bm A.
\]
定义中$\bm A$与$\bm A^{-1}$的地位对称：$\bm A^{-1}\bm A=\bm A\bm A^{-1}=\bm I_n$同样说明$\bm A$是$\bm A^{-1}$的逆。

\paragraph{数乘的逆}
\[
 (c\bm A)^{-1}=c^{-1}\bm A^{-1}\qquad(c\neq0).
\]
直接相乘：$(c\bm A)(c^{-1}\bm A^{-1})=(cc^{-1})\bm A\bm A^{-1}=\bm I_n$，反方向同理。

\paragraph{转置的逆}
\[
 (\bm A^\top)^{-1}=(\bm A^{-1})^\top.
\]
由乘积的转置，$\bm A^\top(\bm A^{-1})^\top=(\bm A^{-1}\bm A)^\top=\bm I^\top=\bm I$，反方向同理，所以$(\bm A^{-1})^\top$就是$\bm A^\top$的逆。

\paragraph{正交矩阵的逆}
\[
 \bm Q^{-1}=\bm Q^\top,\qquad\bm Q\bm Q^\top=\bm I_n\qquad(\bm Q\text{为正交矩阵}).
\]
正交矩阵的转置条件$\bm Q^\top\bm Q=\bm I_n$只说明$\bm Q^\top$是$\bm Q$的左逆。方阵只要有一侧的逆就可逆，而且这个单侧逆就是逆矩阵；这一点的证明要用到行列式，放在第~\ref{sec:ls-square-one-sided-inverse}节。由此$\bm Q^{-1}=\bm Q^\top$，另一侧的乘积$\bm Q\bm Q^\top$也等于$\bm I_n$。

$\bm Q\bm Q^\top$的$(i,j)$元是$\bm Q$的第$i$行与第$j$行的内积，所以正交矩阵的各行同样两两正交且长度为$1$，求逆只需转置。又因$(\bm Q^\top)^\top\bm Q^\top=\bm Q\bm Q^\top=\bm I_n$，由转置条件，$\bm Q^{-1}=\bm Q^\top$也是正交矩阵。

\subsubsection{幂的性质}

设$\bm A$为方阵，$m,n$为非负整数；$\bm A$可逆时，$m,n$可取任意整数。

\paragraph{同底幂相乘}
\[
 \bm A^m\bm A^n=\bm A^{m+n}.
\]
由结合律，两边都是同样个数的$\bm A$（或$\bm A^{-1}$）连乘；指数异号时，相邻的$\bm A\bm A^{-1}$逐对消去。交换$m,n$的角色可见，同一矩阵的不同次幂彼此可以交换。

\paragraph{幂的幂}
\[
 (\bm A^m)^n=\bm A^{mn}.
\]
左边是$n$组、每组$m$个因子的连乘，由结合律去掉括号即得。

\paragraph{负幂与逆相容}
\[
 (\bm A^k)^{-1}=(\bm A^{-1})^k=\bm A^{-k}.
\]
对$k$个相同因子反复使用式~\eqref{eq:ls-inverse-product}即得。

\subsubsection{迹的性质}
\phantomsection
\label{sec:ls-matrix-contractions}

\paragraph{线性}
\[
 \operatorname{tr}(a\bm A+b\bm B)=a\operatorname{tr}(\bm A)+b\operatorname{tr}(\bm B),
\]
这是有限和的线性。

\paragraph{转置不变}
\[
 \operatorname{tr}(\bm A^\top)=\operatorname{tr}(\bm A),
\]
因为转置不改变对角元。

\paragraph{循环性}
对$\bm A\in\mathbb R^{m\times n}$、$\bm B\in\mathbb R^{n\times m}$，
\begin{equation}
 \operatorname{tr}(\bm A\bm B)
 =\sum_{i=1}^m\sum_{j=1}^n A_{i,j}B_{j,i}
 =\operatorname{tr}(\bm B\bm A).
 \label{eq:linear-algebra-trace-cyclic}
\end{equation}
两个双重和只是求和次序不同。$\bm A\bm B$与$\bm B\bm A$可以是不同阶的方阵，相等的是迹，不是矩阵本身，下面的例子说明这一点。
\begin{example}[不同阶乘积的迹相等]
\label{ex:ls-trace-cyclic-different-sizes}
给定：$\bm A=\left[\begin{smallmatrix}1&2&3\end{smallmatrix}\right]\in\mathbb R^{1\times3}$，$\bm B=(1,0,2)^\top\in\mathbb R^{3\times1}$。
计算：$\bm A\bm B=1\cdot1+2\cdot0+3\cdot2=7$是$1\times1$矩阵；
\[
 \bm B\bm A=\begin{bmatrix}1&2&3\\0&0&0\\2&4&6\end{bmatrix}
\]
是$3\times3$矩阵，对角元为$1,0,6$。
结论：$\operatorname{tr}(\bm A\bm B)=7=1+0+6=\operatorname{tr}(\bm B\bm A)$，而两个乘积的阶数不同。
\end{example}

三个因子时，循环性只允许把因子整体循环移动：
\[
 \operatorname{tr}(\bm A\bm B\bm C)=\operatorname{tr}(\bm B\bm C\bm A)=\operatorname{tr}(\bm C\bm A\bm B).
\]
第一个等号把$\bm A$与$\bm B\bm C$看作两个因子使用式~\eqref{eq:linear-algebra-trace-cyclic}，第二个等号把$\bm A\bm B$与$\bm C$看作两个因子。任意交换相邻因子则可能改变迹，下面是一个反例。
\begin{example}[反例：交换相邻因子会改变迹]
\label{ex:ls-trace-adjacent-swap}
给定：$\bm A=\left(\begin{smallmatrix}1&0\\0&0\end{smallmatrix}\right)$，$\bm B=\left(\begin{smallmatrix}0&0\\1&0\end{smallmatrix}\right)$，$\bm C=\left(\begin{smallmatrix}0&1\\0&0\end{smallmatrix}\right)$。
计算：$\bm A\bm B=\bm0$，所以$\bm A\bm B\bm C=\bm0$。另一方面$\bm A\bm C=\bm C$，所以
\[
 \bm A\bm C\bm B=\begin{pmatrix}0&1\\0&0\end{pmatrix}\begin{pmatrix}0&0\\1&0\end{pmatrix}=\begin{pmatrix}1&0\\0&0\end{pmatrix}.
\]
结论：$\operatorname{tr}(\bm A\bm B\bm C)=0$，而$\operatorname{tr}(\bm A\bm C\bm B)=1$；交换$\bm B$与$\bm C$的次序改变了迹。
\end{example}

\paragraph{Frobenius内积的迹表示}
\[
 \langle\bm A,\bm B\rangle_{\mathrm F}=\operatorname{tr}(\bm A^\top\bm B),\qquad
 \lVert\bm A\rVert_{\mathrm F}^2=\operatorname{tr}(\bm A^\top\bm A).
\]
$\bm A^\top\bm B$的第$j$个对角元为$\sum_i A_{i,j}B_{i,j}$，再对$j$求和恰好遍历所有位置。反过来，计算$\operatorname{tr}(\bm A\bm B)$时也不必形成乘积：只需算$\sum_{i,j}A_{i,j}B_{j,i}$，代价为$mn$次乘法，而形成整个$\bm A\bm B$需要$m^2n$次。

\subsubsection{行列式的性质}
\label{sec:ls-determinant-properties}
\phantomsection
\label{sec:prelim-determinants}

行列式由楔积的两条规则定义，下面前两条性质就是这两条规则在列上的表述，其余性质由它们推出。设$\bm A,\bm B$为$d$阶方阵。

\paragraph{对列的多线性}
\[
 \det\begin{bmatrix}\cdots&a\bm u+b\bm w&\cdots\end{bmatrix}
 =a\det\begin{bmatrix}\cdots&\bm u&\cdots\end{bmatrix}
 +b\det\begin{bmatrix}\cdots&\bm w&\cdots\end{bmatrix},
\]
其中三个矩阵只在同一列不同，其余各列相同。$\det\bm A$是各列楔积的系数，所以这正是楔积对每个因子保持线性组合的直接翻译。

\paragraph{交替性}
交换两列，行列式变号；两列相同，行列式为零。前者来自楔积交换两个因子时变号（不相邻的两个因子可以经奇数次相邻交换互换位置）；后者是因为交换两个相同的列不改变矩阵，行列式却要变号，只能为零。

两条合起来又得到两个常用推论。把某列的倍数加到另一列上，行列式不变：按多线性拆开，多出的一项含两列相同，为零。某一列是其余各列的线性组合时，行列式为零：把这一列按组合展开，每一项都含两列相同。例如$\det\left(\begin{smallmatrix}1&2\\2&4\end{smallmatrix}\right)=4-4=0$，第二列是第一列的$2$倍。

\paragraph{转置不变}
\[
 \det(\bm A^\top)=\det\bm A.
\]
在排列展开中，把每一项的乘积按列号重排，得$\prod_ia_{i,\pi(i)}=\prod_ja_{\pi^{-1}(j),j}$；排列展开的推导已经说明$\operatorname{sgn}(\pi^{-1})=\operatorname{sgn}(\pi)$。$\pi$遍历$S_d$时，$\sigma=\pi^{-1}$也遍历$S_d$，所以
\[
 \det\bm A=\sum_{\sigma\in S_d}\operatorname{sgn}(\sigma)\prod_{j=1}^da_{\sigma(j),j}
 =\sum_{\sigma\in S_d}\operatorname{sgn}(\sigma)\prod_{j=1}^d(\bm A^\top)_{j,\sigma(j)}
 =\det(\bm A^\top).
\]
于是对列成立的多线性与交替性，对行同样成立。

\paragraph{数乘与加法}
\[
 \det(c\bm A)=c^d\det\bm A,\qquad
 \det(\bm A+\bm B)\neq\det\bm A+\det\bm B\ (\text{一般}).
\]
数乘把$d$列各乘以$c$，由对每列的线性提出$d$个$c$。行列式只对每一列分别线性，对整个矩阵却不是线性的，下面是一个反例。
\begin{example}[反例：行列式不保持加法]
\label{ex:ls-det-not-additive}
给定：$\bm A=\bm B=\bm I_2$。
计算：$\det\bm A+\det\bm B=1+1=2$，而由第一式，$\det(\bm A+\bm B)=\det(2\bm I_2)=2^2\det\bm I_2=4$。
结论：$\det(\bm A+\bm B)\neq\det\bm A+\det\bm B$。
\end{example}

\paragraph{乘法性}
\[
 \det(\bm A\bm B)=\det\bm A\,\det\bm B.
\]
接连作用两次时，体积缩放倍数应当相乘。排列展开的推导只用到楔积的两条规则，所以把$\bm e_1,\ldots,\bm e_d$换成任意向量$\bm w_1,\ldots,\bm w_d$，同样有$\bigwedge_j\bigl(\sum_ib_{i,j}\bm w_i\bigr)=(\det\bm B)\,\bm w_1\wedge\cdots\wedge\bm w_d$。由$(\bm A\bm B)\bm e_j=\bm A(\bm B\bm e_j)=\sum_i b_{i,j}\bm A\bm e_i$，取$\bm w_i=\bm A\bm e_i$，得到
\[
 (\bm A\bm B)\bm e_1\wedge\cdots\wedge(\bm A\bm B)\bm e_d
 =(\det\bm B)\,\bm A\bm e_1\wedge\cdots\wedge\bm A\bm e_d
 =(\det\bm A)(\det\bm B)\,\bm e_1\wedge\cdots\wedge\bm e_d.
\]
尽管矩阵乘法不可交换，$\det(\bm A\bm B)=\det(\bm B\bm A)$总成立，因为两边都等于$\det\bm A\,\det\bm B$。

\paragraph{三角矩阵的行列式}
\[
 \det\bm A=a_{1,1}a_{2,2}\cdots a_{d,d}\qquad(\bm A\text{为上三角或下三角矩阵}).
\]
设$\bm A$为上三角矩阵。排列$\pi$不是恒等排列时，必有某行$i$满足$\pi(i)<i$（因为$\sum_i(\pi(i)-i)=0$而各项不全为零），这一行选到主对角线下方的零元素，所以这一项为零；排列展开中只剩恒等排列一项。下三角时同理。例如$\det\left(\begin{smallmatrix}1&2&3\\0&4&5\\0&0&6\end{smallmatrix}\right)=1\cdot4\cdot6=24$；特别地，$\det\operatorname{diag}(d_1,\ldots,d_d)=d_1\cdots d_d$，$\det\bm I=1$。与乘法性合起来，两个上三角矩阵乘积的行列式就是两者全部对角元之积，这与三角矩阵乘积的对角元逐个相乘一致。

\paragraph{按行（列）展开}
对$d\geq2$，删除第$i$行、第$j$列得到的$(d-1)$阶子矩阵记为$\bm A^{(i,j)}$，定义\term{代数余子式}（cofactor）
\[
 C_{i,j}=(-1)^{i+j}\det\bm A^{(i,j)}.
\]
则对任一行$i$与任一列$j$，
\[
 \det\bm A=\sum_{j=1}^d a_{i,j}C_{i,j}=\sum_{i=1}^d a_{i,j}C_{i,j}.
\]
把排列展开按第$i$行选取的列分类，选第$j$列的各项之和正是$a_{i,j}C_{i,j}$；按列同理。下面沿第一行展开示例~\ref{ex:ls-det-three-by-three}中的矩阵，与排列展开的结果对照。
\begin{example}[沿第一行展开三阶行列式]
\label{ex:ls-det-cofactor-expansion}
给定：$\bm A=\left(\begin{smallmatrix}2&0&1\\1&3&0\\0&1&1\end{smallmatrix}\right)$。
计算：第一行元素为$2,0,1$，对应代数余子式的符号依次为$+,-,+$，
\[
 \begin{aligned}
 \det\bm A
 &=2\cdot\det\begin{pmatrix}3&0\\1&1\end{pmatrix}
 -0\cdot\det\begin{pmatrix}1&0\\0&1\end{pmatrix}
 +1\cdot\det\begin{pmatrix}1&3\\0&1\end{pmatrix}\\
 &=2\cdot3-0\cdot1+1\cdot1=7.
 \end{aligned}
\]
结论：结果与示例~\ref{ex:ls-det-three-by-three}的排列展开一致。
\end{example}

\paragraph{古典伴随恒等式}
把全部代数余子式排成矩阵$\bm C$，其转置称为\term{古典伴随}（classical adjugate）$\operatorname{adj}(\bm A)=\bm C^\top$。
\begin{theorem}[古典伴随恒等式]
\label{thm:ls-adjugate-identity}
对任意$d$阶方阵$\bm A$，
\[
 \bm A\,\operatorname{adj}(\bm A)
 =\operatorname{adj}(\bm A)\,\bm A
 =(\det\bm A)\,\bm I_d.
\]
\end{theorem}
\begin{proof}
乘积$\bm A\bm C^\top$的$(i,k)$元素为$\sum_j a_{i,j}C_{k,j}$。当$i=k$时，这正是沿第$i$行展开的$\det\bm A$。当$i\neq k$时，$C_{k,j}$不涉及第$k$行的元素，所以这个和是把第$k$行换成第$i$行后的矩阵沿第$k$行的展开；新矩阵有两行相同，行列式为零。于是$\bm A\bm C^\top=(\det\bm A)\bm I_d$。按列展开同样得到$\bm C^\top\bm A=(\det\bm A)\bm I_d$。
\end{proof}
二阶时$\operatorname{adj}\left(\begin{smallmatrix}a&b\\c&d\end{smallmatrix}\right)=\left(\begin{smallmatrix}d&-b\\-c&a\end{smallmatrix}\right)$，恒等式两边除以$ad-bc$就是上一小节的二阶求逆公式。

古典伴随是由余子式得到的代数构造，与第~\ref{chap:linear-representation}章由内积配对定义的伴随不是同一个对象。高阶时展开余子式的代价是阶乘级的，所以恒等式主要作为证明工具：下面用它证明方阵的单侧逆即逆，第~\ref{sec:ls-determinant-invertibility}节用它给出可逆判据，第~\ref{sec:ls-characteristic-polynomial}节用它证明Cayley--Hamilton定理。

\paragraph{方阵的单侧逆即逆}
\phantomsection
\label{sec:ls-square-one-sided-inverse}
若$n$阶方阵满足$\bm B\bm A=\bm I_n$或$\bm A\bm B=\bm I_n$，则$\bm A$可逆且$\bm B=\bm A^{-1}$。

由乘法性，两种情形都有$\det\bm A\det\bm B=\det\bm I_n=1$，所以$\det\bm A\neq0$。令$\bm C=\operatorname{adj}(\bm A)/\det\bm A$，古典伴随恒等式给出$\bm A\bm C=\bm C\bm A=\bm I_n$，所以$\bm A$可逆，$\bm C=\bm A^{-1}$。若$\bm B\bm A=\bm I_n$，则$\bm B=\bm B(\bm A\bm C)=(\bm B\bm A)\bm C=\bm C$；若$\bm A\bm B=\bm I_n$，则$\bm B=(\bm C\bm A)\bm B=\bm C(\bm A\bm B)=\bm C$。所以检查方阵的逆只需验证一侧的乘积。对矩形矩阵，左逆与右逆则可以只存在一个，也可以不唯一（示例~\ref{ex:ls-left-inverse-not-unique}）。

\paragraph{正交矩阵的行列式}
\[
 \det\bm Q=\pm1\qquad(\bm Q\text{为正交矩阵}).
\]
由$\det(\bm Q^\top)\det\bm Q=\det(\bm Q^\top\bm Q)=\det\bm I_n=1$及$\det(\bm Q^\top)=\det\bm Q$，得$(\det\bm Q)^2=1$。旋转$\bm R_\phi$的行列式为$\cos^2\phi+\sin^2\phi=1$；交换两个坐标的置换矩阵行列式为$-1$，它翻转了定向。

\subsubsection{范数的性质}

\paragraph{列和与行和公式}
\[
 \lVert\bm A\rVert_1=\max_j\sum_i|a_{i,j}|,\qquad
 \lVert\bm A\rVert_\infty=\max_i\sum_j|a_{i,j}|.
\]
第一式按列算：
\[
 \lVert\bm A\bm x\rVert_1=\Bigl\lVert\sum_jx_j\bm A_{:,j}\Bigr\rVert_1\leq\sum_j|x_j|\,\lVert\bm A_{:,j}\rVert_1\leq\Bigl(\max_j\lVert\bm A_{:,j}\rVert_1\Bigr)\lVert\bm x\rVert_1,
\]
而取$\bm x$为最大绝对列和所在列对应的标准坐标向量时等号成立。第二式按行算：对每个$i$，
\[
 |(\bm A\bm x)_i|\leq\sum_j|a_{i,j}|\,|x_j|\leq\Bigl(\sum_j|a_{i,j}|\Bigr)\lVert\bm x\rVert_\infty,
\]
对$i$取最大值即得上界；取输入各坐标为最大绝对行和那一行元素的符号（元素为零处取$1$），等号成立。

\paragraph{谱范数不超过Frobenius范数}
\[
 \lVert\bm A\bm x\rVert_2\leq\lVert\bm A\rVert_{\mathrm F}\lVert\bm x\rVert_2,\qquad
 \lVert\bm A\rVert_2\leq\lVert\bm A\rVert_{\mathrm F}.
\]
对每行用$p=q=2$的Hölder不等式，$(\bm A\bm x)_i^2=(\bm A_{i,:}\bm x)^2\leq\lVert\bm A_{i,:}\rVert_2^2\lVert\bm x\rVert_2^2$；对$i$求和，右侧的$\sum_i\lVert\bm A_{i,:}\rVert_2^2$正是$\lVert\bm A\rVert_{\mathrm F}^2$，得到第一式。第一式说明每个比值$\lVert\bm A\bm x\rVert_2/\lVert\bm x\rVert_2$都不超过$\lVert\bm A\rVert_{\mathrm F}$，所以谱范数有限，并得到第二式。

对一般的输入、输出范数，诱导范数同样有限。由第~\ref{sec:ls-error-scale-comparison}节的范数等价，存在常数$c_1,c_2>0$使$\lVert\bm y\rVert_\beta\leq c_1\lVert\bm y\rVert_2$、$\lVert\bm x\rVert_2\leq c_2\lVert\bm x\rVert_\alpha$，于是$\lVert\bm A\bm x\rVert_\beta\leq c_1c_2\lVert\bm A\rVert_{\mathrm F}\lVert\bm x\rVert_\alpha$。诱导范数还满足范数的四条要求：非负性与齐次性由定义直接得到；若$\lVert\bm A\rVert_{\alpha\to\beta}=0$，则每个$\bm A\bm e_j=\bm0$，即$\bm A=\bm0$；三角不等式来自$\lVert(\bm A+\bm B)\bm x\rVert_\beta\leq\lVert\bm A\bm x\rVert_\beta+\lVert\bm B\bm x\rVert_\beta$，两边除以$\lVert\bm x\rVert_\alpha$再取上确界即得。此外，定义中的上确界在某个非零方向上达到，因而也可以写成最大值；这一点要用范数等价和逐坐标抽取收敛子列的论证，这里不展开。

\paragraph{次乘性}
中间范数匹配时，
\[
 \lVert\bm A\bm B\rVert_{\alpha\to\gamma}
 \leq\lVert\bm A\rVert_{\beta\to\gamma}
 \lVert\bm B\rVert_{\alpha\to\beta}.
\]
连续应用两次诱导范数的界：
\[
 \lVert\bm A\bm B\bm x\rVert_\gamma\leq\lVert\bm A\rVert_{\beta\to\gamma}\lVert\bm B\bm x\rVert_\beta\leq\lVert\bm A\rVert_{\beta\to\gamma}\lVert\bm B\rVert_{\alpha\to\beta}\lVert\bm x\rVert_\alpha.
\]
两边除以$\lVert\bm x\rVert_\alpha$再取上确界即得。特别地，反复使用可得$\lVert\bm A^k\rVert_p\leq\lVert\bm A\rVert_p^k$。

\paragraph{Frobenius范数与谱范数的混合界}
\begin{equation}
 \lVert\bm A\bm B\rVert_{\mathrm F}
 \leq\lVert\bm A\rVert_2\lVert\bm B\rVert_{\mathrm F},
 \qquad
 \lVert\bm A\bm B\rVert_{\mathrm F}
 \leq\lVert\bm A\rVert_{\mathrm F}\lVert\bm B\rVert_2.
 \label{eq:matrix-analysis-mixed-submultiplicativity}
\end{equation}
第一式按列算：
\[
 \lVert\bm A\bm B\rVert_{\mathrm F}^2=\sum_j\lVert\bm A\bm B_{:,j}\rVert_2^2\leq\lVert\bm A\rVert_2^2\sum_j\lVert\bm B_{:,j}\rVert_2^2=\lVert\bm A\rVert_2^2\lVert\bm B\rVert_{\mathrm F}^2.
\]
第二式对转置应用第一式：$\lVert\bm A\bm B\rVert_{\mathrm F}=\lVert\bm B^\top\bm A^\top\rVert_{\mathrm F}\leq\lVert\bm B^\top\rVert_2\lVert\bm A^\top\rVert_{\mathrm F}$，其中转置不改变Frobenius范数，还需要说明$\lVert\bm B^\top\rVert_2=\lVert\bm B\rVert_2$。对固定的$\bm x$，由$p=q=2$的Hölder不等式，$\lVert\bm y\rVert_2=1$时$|\bm y^\top\bm B\bm x|\leq\lVert\bm B\bm x\rVert_2$，而$\bm B\bm x\neq\bm0$时取$\bm y=\bm B\bm x/\lVert\bm B\bm x\rVert_2$即取等号，所以
\[
 \lVert\bm B\rVert_2=\sup_{\lVert\bm x\rVert_2=\lVert\bm y\rVert_2=1}|\bm y^\top\bm B\bm x|.
\]
由于$\bm y^\top\bm B\bm x=\bm x^\top\bm B^\top\bm y$，交换$\bm x,\bm y$的角色，同一个上确界也等于$\lVert\bm B^\top\rVert_2$。结合$\lVert\bm A\rVert_2\leq\lVert\bm A\rVert_{\mathrm F}$，Frobenius范数也满足次乘性：$\lVert\bm A\bm B\rVert_{\mathrm F}\leq\lVert\bm A\rVert_{\mathrm F}\lVert\bm B\rVert_{\mathrm F}$。

\paragraph{正交不变性}
设$\bm Q$为正交矩阵，则
\[
 \lVert\bm Q\bm x\rVert_2=\lVert\bm x\rVert_2,\qquad
 \lVert\bm Q\rVert_2=1,
\]
\[
 \lVert\bm Q\bm A\rVert_{\mathrm F}=\lVert\bm A\rVert_{\mathrm F},\qquad
 \lVert\bm Q\bm A\rVert_2=\lVert\bm A\rVert_2.
\]
由$\lVert\bm Q\bm x\rVert_2^2=\bm x^\top\bm Q^\top\bm Q\bm x=\lVert\bm x\rVert_2^2$，正交矩阵保持欧氏长度，所以每个比值$\lVert\bm Q\bm x\rVert_2/\lVert\bm x\rVert_2$都等于$1$，$\lVert\bm Q\rVert_2=1$。对矩阵$\bm A$，左乘$\bm Q$不改变各列的长度，$\lVert\bm Q\bm A_{:,j}\rVert_2=\lVert\bm A_{:,j}\rVert_2$，对$j$求平方和即得$\lVert\bm Q\bm A\rVert_{\mathrm F}=\lVert\bm A\rVert_{\mathrm F}$；它也不改变每个输出的长度，$\lVert\bm Q\bm A\bm x\rVert_2=\lVert\bm A\bm x\rVert_2$，取上确界即得$\lVert\bm Q\bm A\rVert_2=\lVert\bm A\rVert_2$。
